\section{The 2025 filter and the linear guide}
\label{sec:filter}

Consider the filter of \citet{whitaker2025}. Let \(X_t\) denote prevalence and let \(N_t\) denote cumulative incidence. The unconditioned hazard is \(h\), and \(H = \mathrm{diag}\{h\}\). A scalar observation \(y_{t_i}\) is a noisy count of one incidence component over \((t_{i-1}, t_i]\), selected by a vector \(P\). In the susceptible-infected-removed applications, \(P' = (1, 0)\). The contact rate enters on the log scale as \(\tilde{\beta}_t\), in the precision form printed as equation (2) of \citet{whitaker2025}. The static parameter and the particle index are those of their Algorithm 2. The comparison below holds the rate parameters fixed, so the gamma updates of their equations (9)--(11) stay outside the calculation.

Their Algorithm 2 is printed for the negative binomial observation. The synthetic study in their Section 4.1 uses the binomial observation, and the observation factor in the weight is the binomial probability mass function.

\begin{algbox}
\label{alg:schematic}
Schematic of Algorithm 2 of \citet{whitaker2025}. The draw from \(q\) in the propagation step is the linear incidence guide. Section~\ref{sec:residual} replaces the Poisson means of that draw.
\begin{enumerate}
\item Propagate each particle. Jitter the observation parameter with the kernel printed for the negative binomial model. Draw the contact-rate path from its transition, which the paper calls a blind propagation. Draw the incidence path from the proposal \(q\), which depends on the next observation.
\item Resample with the weight in Equation~\eqref{eq:weight}.
\item Update the sufficient statistic for the rate parameters.
\item Draw those rate parameters from their conditional posteriors.
\end{enumerate}
\end{algbox}

The weight in step 2, for particle \(k\) on the window that ends at the next observation, is
\begin{equation}
\label{eq:weight}
w^{(k)} \propto
\frac{\pi\bigl(\Delta n^{(k)} \mid \tilde{\beta}^{(k)}, \theta^{(k)}\bigr)\,
\pi\bigl(y \mid \Delta n^{(k)}, \phi^{(k)}\bigr)}
{q\bigl(\Delta n^{(k)} \mid \tilde{\beta}^{(k)}, \phi^{(k)}, \theta^{(k)}, y\bigr)}.
\end{equation}
The numerator is the time-discretised transition and the observation mass function. It does not contain the conditioned hazard \(h^{*}\). The denominator is \(q\). We use the ratio as an importance weight for one propagate step.

The proposal \(q\) is built in their Section 3.4. At a time \(\tau\) inside the observation window, equation (12) of \citet{whitaker2025} approximates the remaining incidence by a Gaussian whose mean and variance are the current hazard times the remaining time. The observation is given a Gaussian approximation, and the mean inside the observation variance is replaced by the mean of that Gaussian. The display after their equation (18) is the conditioned hazard \(h^{*}\). End-point paths are their forward recursion with \(h\) replaced by \(h^{*}\). The density \(q\) is the product of Poisson probabilities with rates \(h^{*}\Delta\tau\).

The printed correction multiplies by \(\rho P' H\), a row when the observation is scalar. The cross-covariance block that enters the conditional mean of the remaining incidence is a column. For a diagonal \(H\) the row and the column have the same entries. On the susceptible-infected-removed model the shared entries correct the infection component only,
\begin{equation}
\label{eq:hstar}
h^{*}_{\mathrm{inf}}
= h_{\mathrm{inf}}
+ \rho h_{\mathrm{inf}}
\frac{y - \hat{\mu}}{\rho^{2} h_{\mathrm{inf}}(t_i - \tau) + \sigma^{2}},
\qquad
h^{*}_{\mathrm{rem}} = h_{\mathrm{rem}},
\end{equation}
where \(\hat{\mu} = \rho\bigl(n^{\mathrm{past}}_{\mathrm{inf}} + h_{\mathrm{inf}}(t_i - \tau)\bigr)\) and, for the binomial variance evaluated at \(\hat{\mu}\), \(\sigma^{2} = (1-\rho)\hat{\mu}\). Equation~\eqref{eq:hstar} is the scalar form used in the simulation. Its entries match the printed display, and they match equation (9) of \citet{golightly2015} when the stoichiometry is the identity, so that the correction is the column \(HP\).

The study in Section~\ref{sec:simulation} holds the contact-rate path and the prevalence at the left endpoint of the window at the simulated values, and it holds the removal rate and the reporting probability at the values used to simulate the data. The jitter, the blind contact-rate draw, the sufficient-statistic update and the conditional parameter draw are omitted. A particle in Algorithm~\ref{alg:schematic} is propagated without the latent path. The comparison concerns \(h^{*}\) given that path.
